% Shared exercise chunk: l4-fitting-support
\begin{exercisechunk}
\item \textbf{Likelihood accuracy, data, and support.\suppmark}
\label[exercise]{ex:fitting-support}
Parts (a)--(c) use the repeated-error generators from
\cref{ex:binary-count} and all-zero training sequences.
\begin{enumerate}[(a),beginpenalty=10000]
\item Compute both empirical per-token losses. Include
\(P=\delta_{0^T}\) in the model class and show that
\(\varepsilon=1-e^{-\eta}\) makes both generators acceptable at
optimization tolerance \(\eta>0\).
\item Compute their terminal and counting costs and their leading
terms as \(T\eta\to0\). Give a sufficient order of \(\eta\) in
\(T\) to keep each expected count bounded.
\item Remove \(P\) and use a finite collection of paired generators
at distinct positive values of \(\varepsilon\). Identify all
empirical minimizers. Explain why exact fitting leaves a tie between
models of different task cost.
\item Specialize \cref{thm:finite-mle} to zero-cost data, first for
\(c\in[0,T]\) and then for \(c\in[0,1]\). Express each guarantee
using the number of sequences \(n\), then using the token budget
\(L=nT\).
\item In \cref{prop:conditional-support}, choose a prompt on which
the two candidate rules disagree. Show that its expert label identifies
the correct rule. Compare the number of such queries with the number
of random demonstrations needed for failure probability at most
\(\delta\in(0,1/2)\).
\end{enumerate}
\end{exercisechunk}
